\documentclass[12pt]{article}

\usepackage{amsmath, amssymb, amsthm}
\usepackage{enumerate}
\usepackage{hyperref}
\usepackage{geometry}
\usepackage{newtxtext}
\usepackage{newtxmath}
\geometry{margin=1in}

\newtheorem*{claim}{Claim}
\newtheorem*{lemma}{Lemma}

\title{Additive Combinatorics - Lecture 2}
\author{}
\date{}

% ----------------------------------------------------------------
\begin{document}
\maketitle

I want to reproduce Dvir's proof of the finite field Kakeya conjecture, following Terence Tao's write up at his blog.
This is to be done as a self contained write-up, and is made as simple as possible for even a beginner to understand, and to even attempt the proofs of the main statements.

If $F^n$ is a vector space over a finite field $F$,
a \textit{Kakeya set} in $F^n$ is a subset of $F^n$
that contains a line in every direction.

The finite field Kakeya conjecture states that
any Kakeya set in $F^n$ must have size at least $C_n |F|^n$ for some constant $C_n$ depending only on $n$.

% ================================================================
\section*{Questions}

\begin{enumerate}[(a)]

  \item Let $F$ be a field and let $d \geq 0$ be an integer. Let $F[X]$
  denote the polynomial ring in one variable over $F$.
  Prove, using strong induction on $d$, that if $P \in F[X]$ is a
  non-zero polynomial of degree at most $d$, then $P$ has at most $d$
  roots in $F$.
  For the base case take $d = 0$.
  For the inductive step, assume the result holds for all degrees
  $0, 1, \ldots, d-1$ and deduce it for degree $d$.

  \medskip
  \noindent(Hint: a field is an integral domain: the product of any two
  non-zero elements is non-zero.

  \medskip
  \noindent Recall the \textbf{Factor Theorem}: if $r \in F$ is a root
  of $P \in F[X]$, then $P(X) = (X - r)Q(X)$ for some $Q \in F[X]$
  of degree $\deg P - 1$.  You may use this without proof.)

  \medskip
  \noindent\textit{Warm-up.} Let $P(X) = X^2 - 5X + 6$.
  \begin{enumerate}[(i)]
    \item Factorise $P$. How many roots does $P$ have? How does this compare to its degree?
    \item Now suppose $Q(X)$ is \emph{any} non-zero polynomial of degree
    $2$. Explain why $Q$ can have at most $2$ roots.
  \end{enumerate}

  \item Given any finite set $E \subseteq F$ with $|E| \leq d$, show
  that there exists a non-zero polynomial $P \in F[X]$ of degree at
  most $d$ such that $P(x) = 0$ for all $x \in E$.

  \medskip
  \noindent\textit{Warm-up.} Let $F = \mathbb{R}$ and $E = \{1, 3, 7\}$.
  \begin{enumerate}[(i)]
    \item Write down a polynomial of degree $3$ that vanishes at every point of $E$.
    \item Can you write down a non-zero polynomial of degree $2$ that
    vanishes at all three points?
  \end{enumerate}

  \item Let $E \subseteq F^n$ be a set of cardinality less than $\binom{n+d}{d}$.
  Show that there exists a non-zero polynomial
  $P \in F[X_1, \ldots, X_n]$ of degree at most $d$ such that $P(x) = 0$
  for all $x \in E$.

  \medskip
  \noindent\textit{Warm-up.}

  \textbf{Stars and Bars.}
  The stars and bars method counts the number of solutions in $\mathbb{Z}_{\ge 0}$ to
  \[
  x_1 + \cdots + x_k = N.
  \]
  It does this by placing $N$ identical ``stars'' in a row and inserting $k-1$ ``bars'' to split them into $k$ groups. Each arrangement corresponds uniquely to a solution $(x_1,\ldots,x_k)$, so the number of solutions is
  \[
  \binom{N+k-1}{k-1}.
  \]

  \textbf{Kernel.}
  Let $T: V \to W$ be a linear map of vector spaces over $F$. The kernel of $T$ is
  \[
  \ker(T) = \{ v \in V \mid T(v) = 0 \}.
  \]
  It is a subspace of $V$, and $T$ is injective if and only if $\ker(T)=\{0\}$.

  \begin{enumerate}[(i)]
    \item A monomial has the form
    \[
    X_1^{\alpha_1} \cdots X_n^{\alpha_n}, \quad \alpha_i \ge 0.
    \]
    Compute the number of monomials in $n$ variables of total degree at most $d$.
    (Hint: count the number of non-negative integer solutions of
    \[
    \alpha_1 + \cdots + \alpha_n \le d.
    \]
    You may use the stars-and-bars method.)
    \item Let $F^E$ be the vector space of functions $f : E \to F$.
    Compute $\dim_F(F^E)$, and use the hypothesis on $|E|$ to compare it
    with $\dim_F(V)$.

    \item Let $V$ be the vector space of polynomials in $F[X_1, \ldots, X_n]$ of degree at most $d$. Define the evaluation map
    \[
    \mathrm{ev} : V \to F^E, \quad P \mapsto (P(e))_{e \in E}.
    \]
    Compare $\dim_F(V)$ and $\dim_F(F^E)$ and deduce that $\mathrm{ev}$ can not be injective and has a nontrivial kernel.
    \item Using the definition of a nontrivial kernel and the map in the previous exercise is noninjective, conclude that there exists a non-zero polynomial $P \in V$ such that
    \[
    P(e) = 0 \quad \text{for all } e \in E.
    \]
  \end{enumerate}

  Here are the statements of the main theorems.

  \item Let $E$ be a Kakeya set in $F^n$, and let
  $P \in F[X_1, \ldots, X_n]$ be a polynomial of degree at most $|F|-1$
  that vanishes on a Kakeya set $E$.  Show that $P = 0$.
  (Could not find good warm ups for these two remaining exercises)

  \item Show that every Kakeya set in $F^n$ has size at least $\binom{n+|F|-1}{n}$.
  Conclude that every Kakeya set in $F^n$ has size at least
  $C_n |F|^n$ for some constant $C_n$ depending only on $n$.

\end{enumerate}

% ================================================================
\newpage
\section*{Solutions}

Credit: Tao's exposition of Dvir's proof of the finite field Kakeya conjecture,
\href{https://terrytao.wordpress.com/2008/03/24/dvirs-proof-of-the-finite-field-kakeya-conjecture/}{\textit{Tao's proof of the finite field Kakeya conjecture}}.

\begin{enumerate}[(a)]

  \item

  The factorisation of $P(X) = (X - 2)(X - 3)$ admits two roots, $2$
  and $3$, which matches the degree of $P$.

  For any non-zero polynomial $Q$ of degree $2$, if $r$ is a root of
  $Q$, then by the Factor Theorem $Q(X) = (X - r)R(X)$ for some
  $R \in F[X]$ of degree $1$.  If $Q$ had a second root $s \neq r$,
  then $R(s) = 0$ and again by the Factor Theorem
  $R(X) = (X - s)S(X)$ for some non-zero constant $S$.
  Hence $Q(X) = (X - r)(X - s)S(X)$, and $Q$ has no more than $2$ roots.

  For the main proposition, we argue by strong induction on $d$.

  \textit{Base case} $d = 0$: a non-zero constant polynomial has no
  roots at all, so it has at most $0$ roots. \checkmark

  \textit{Inductive step}: Suppose the statement holds for all degrees
  $0, 1, \ldots, d-1$, and let $P$ have degree exactly $d \ge 1$.
  If $P$ has no roots in $F$, the claim holds trivially.
  Otherwise pick a root $r \in F$.
  By the Factor Theorem, $P(X) = (X - r)Q(X)$ for some
  $Q \in F[X]$ of degree $d - 1$.
  Since $F$ is an integral domain, any further root $s$ of $P$ satisfies
  $(s - r)Q(s) = 0$; as $s \neq r$ we get $Q(s) = 0$, so $s$ is a
  root of $Q$.
  By the inductive hypothesis, $Q$ has at most $d - 1$ roots, so $P$
  has at most $1 + (d-1) = d$ roots. \qed

  % ---------------------------------------------------------------
  \item

  The polynomial $P(X) = (X-1)(X-3)(X-7)$ has degree $3$ and vanishes
  at every point of $E$.  There is no non-zero polynomial of degree $2$
  that vanishes at all three points: such a polynomial would have at
  least $3 > 2$ roots, contradicting part~(a).

  Let $E = \{e_1, e_2, \ldots, e_k\}$ with $k = |E| \leq d$.

  Construct:
  \[
    P(X) \;=\; \prod_{i=1}^{k}(X - e_i).
  \]
  This polynomial is non-zero (it has leading coefficient $1$),
  has degree exactly $k \leq d$, and satisfies $P(e_i) = 0$ for every
  $e_i \in E$.  \qed

  % ---------------------------------------------------------------
  \item

  The space $V$ of polynomials in $F[X_1, \ldots, X_n]$ of degree at
  most $d$ has as a basis the monomials
  $X_1^{\alpha_1} \cdots X_n^{\alpha_n}$ with
  $\alpha_1 + \cdots + \alpha_n \leq d$.
  Introducing a slack variable $\alpha_{n+1} = d -
  (\alpha_1 + \cdots + \alpha_n) \geq 0$, we count solutions to
  $\alpha_1 + \cdots + \alpha_{n+1} = d$ in non-negative integers.
  By stars and bars, this is $\binom{d+n}{n} = \binom{n+d}{d}$,
  so $\dim_F(V) = \binom{n+d}{d}$.

  The vector space $F^E$ of $F$-valued functions on $E$ has a basis
  consisting of the indicator functions $\mathbf{1}_e$ for $e \in E$
  (where $\mathbf{1}_e(e) = 1$ and $\mathbf{1}_e(e') = 0$ for
  $e' \neq e$), so $\dim_F(F^E) = |E|$.

  By hypothesis, $|E| < \binom{n+d}{d}$, so
  $\dim_F(F^E) = |E| < \binom{n+d}{d} = \dim_F(V)$.

  Since $\dim_F(V) > \dim_F(F^E)$, the evaluation map
  $\mathrm{ev} : V \to F^E$ cannot be injective (a linear map from a
  larger space to a smaller one cannot be injective).
  By definition of the kernel, this means $\ker(\mathrm{ev}) \neq
  \{0\}$: there exists a non-zero polynomial $P \in V$ with
  $\mathrm{ev}(P) = 0$, i.e.\ $P(e) = 0$ for all $e \in E$.  \qed

  % ---------------------------------------------------------------
  \item

  Suppose for contradiction that $P \neq 0$.
  Since $P$ vanishes on the non-empty set $E$ and a non-zero constant
  never vanishes, we have $d := \deg P \geq 1$.
  Write $P = \sum_{i=0}^{d} P_i$ where $P_i$ is the homogeneous
  component of degree $i$; in particular $P_d \neq 0$.

  Let $v \in F^n \setminus \{0\}$ be any direction.
  Since $E$ is a Kakeya set, there exists $x_v \in F^n$ such that the
  line $\{x_v + tv : t \in F\}$ is entirely contained in $E$.
  The map $t \mapsto P(x_v + tv)$ is a polynomial in $t$ of degree at
  most $d \leq |F|-1$.  It vanishes at every $t \in F$ (all $|F|$
  elements), so by part~(a) it is identically zero as a polynomial in
  $t$.

  In particular every coefficient in $t$ is zero.
  The coefficient of $t^d$ in $P(x_v + tv)$ equals $P_d(v)$: indeed,
  the only contribution of degree $d$ in $t$ comes from the leading
  term $P_d(x_v + tv)$, whose $t^d$ coefficient is $P_d(v)$ by
  homogeneity.
  Hence $P_d(v) = 0$ for every $v \in F^n \setminus \{0\}$.
  Since $P_d$ is homogeneous of degree $d \geq 1$, we also have
  $P_d(0) = 0$, so $P_d$ vanishes on all of $F^n$.

  We now derive the contradiction $P_d = 0$ as a polynomial.
  Write
  \[
    P_d(X_1, \ldots, X_n)
    \;=\; \sum_{j=0}^{d} X_1^j \cdot Q_j(X_2, \ldots, X_n)
  \]
  where each $Q_j$ is a polynomial in $X_2, \ldots, X_n$.
  For any fixed $(a_2, \ldots, a_n) \in F^{n-1}$, the single-variable
  polynomial $\sum_{j=0}^{d} Q_j(a_2,\ldots,a_n)\, X_1^j$ has degree
  at most $d \leq |F|-1$ and vanishes at every element of $F$.
  By part~(a), each coefficient is zero: $Q_j(a_2,\ldots,a_n) = 0$
  for all $j$ and all $(a_2,\ldots,a_n) \in F^{n-1}$.
  Applying the same argument to each $Q_j$ in the variable $X_2$
  (with the remaining variables fixed), and continuing variable by
  variable, we conclude that every $Q_j$ is the zero polynomial.
  Hence $P_d = 0$, contradicting $P_d \neq 0$.
  Therefore $P = 0$.  \qed

  % ---------------------------------------------------------------
  \item

  Suppose for contradiction that the Kakeya set $E$ satisfies
  $|E| < \binom{n+|F|-1}{n}$.
  Setting $d = |F|-1$ in part~(c), there exists a non-zero polynomial
  $P \in F[X_1, \ldots, X_n]$ of degree at most $|F|-1$ vanishing on
  $E$.  But by part~(d), any polynomial of degree at most $|F|-1$
  vanishing on a Kakeya set must equal zero.
  This is a contradiction, so $|E| \geq \binom{n+|F|-1}{n}$.

  It remains to show $\binom{n+|F|-1}{n} \geq C_n |F|^n$.
  We have
  \[
    \binom{n+|F|-1}{n}
    \;=\; \frac{|F|\cdot(|F|+1)\cdots(|F|+n-1)}{n!}
    \;\geq\; \frac{|F|^n}{n!},
  \]
  since each of the $n$ factors in the numerator is at least $|F|$.
  Hence every Kakeya set in $F^n$ has size at least $C_n|F|^n$ with
  $C_n = \dfrac{1}{n!}$.  \qed

\end{enumerate}

\end{document}